Safety filters based on control barrier functions (CBFs) minimally modify an unsafe nominal command so that a barrier condition holds~\cite{ames2017control,ames2019control}. The theory is written in continuous time, whereas implementations hold the input constant over a sampling period, and discrete-time CBF conditions have been proposed as an alternative~\cite{zeng2021safety,xiong2023discrete}; sampled-data safety has also been analysed directly, including conditions for forward invariance under piecewise-constant controllers~\cite{breeden2021control}, safety with approximate discrete-time models~\cite{taylor2022safety} and sample-and-hold safety through input-to-state safety~\cite{bahati2023sample}. Practitioners therefore face two knobs, the class-K gain $\alpha$ and the period $\Delta t$, and a trade-off: conservative settings keep the robot far from obstacles but slow it down.

We ask how $\alpha$ and $\Delta t$ affect safety violations and conservatism (time to goal, minimum clearance) for three filters: a continuous-time CBF applied at discrete steps, a discrete-time CBF condition, and a distance-threshold braking heuristic. Both CBF filters impose one scalar constraint and have a closed-form projection (a half-space for the continuous-time condition, the exterior of a ball for the discrete-time one). We state four falsifiable hypotheses (Sec.~\ref{sec:setup}) and report that one (H2) is refuted and three are partly supported; of those three, H1 and H4 fail their stated tests as written and hold only in a weaker form, and H3 holds on the double integrator only. Contributions: (i) a reproducible numpy benchmark of these filters, (ii) a full $\alpha\times\Delta t$ grid showing that one-factor sweeps miss the corner where both are large, with a breakdown of violations into those seen at the control sample instants and those that occur only between them, and (iii) an ablation showing that the obstacle-selection rule for the single constraint matters more than $\Delta t$ for the continuous-time filter at the default gain. This is a small study and makes no claim of a general ranking.
